\documentclass[11pt]{article}

\usepackage[a4paper,margin=1in]{geometry}
\usepackage{amsmath,amssymb,amsthm,mathtools}
\usepackage{mathrsfs}
\usepackage{hyperref}
\usepackage{enumitem}
\usepackage{microtype}
\usepackage{tikz}
\usetikzlibrary{arrows.meta,calc}

\hypersetup{
	colorlinks=true,
	linkcolor=blue,
	urlcolor=blue,
	citecolor=blue
}

\title{Lecture Notes on the Pentagonal Number Theorem}
\author{}
\date{}

\newtheorem{theorem}{Theorem}[section]
\newtheorem{proposition}[theorem]{Proposition}
\newtheorem{lemma}[theorem]{Lemma}
\newtheorem{corollary}[theorem]{Corollary}
\theoremstyle{definition}
\newtheorem{definition}[theorem]{Definition}
\newtheorem{example}[theorem]{Example}
\newtheorem{exercise}[theorem]{Exercise}
\theoremstyle{remark}
\newtheorem{remark}[theorem]{Remark}

\newcommand{\ZZ}{\mathbf{Z}}
\newcommand{\NN}{\mathbf{N}}
\newcommand{\PP}{\mathcal{P}}
\newcommand{\qP}{\mathscr{P}}
\newcommand{\ord}{\operatorname{ord}}
\newcommand{\gen}[1]{\left\langle #1 \right\rangle}

\begin{document}
	
	\maketitle
	
	\tableofcontents
	
	\section{Introduction}
	
	One of the most beautiful identities in the theory of partitions is Euler's \emph{Pentagonal Number Theorem}:
	\[
	\prod_{m=1}^{\infty}(1-q^m)
	=
	\sum_{k=-\infty}^{\infty}(-1)^k q^{k(3k-1)/2}.
	\]
	At first sight this formula is mysterious. The left-hand side is an infinite product, while the right-hand side is a highly sparse power series involving the numbers
	\[
	\frac{k(3k-1)}{2},
	\qquad k\in \ZZ.
	\]
	These are the \emph{generalized pentagonal numbers}:
	\[
	0,1,2,5,7,12,15,22,26,\dots
	\]
	
	The theorem is central in the study of partitions. It leads directly to Euler's recurrence for the partition function \(p(n)\), and it is also a prototype for many deep identities in \(q\)-series, modular forms, and combinatorics.
	
	The aim of these notes is to explain the theorem carefully, from basic definitions up to proof and applications.
	
	\section{Partitions and Generating Functions}
	
	\subsection{Partitions}
	
	\begin{definition}
		A \emph{partition} of a nonnegative integer \(n\) is a way of writing \(n\) as a sum of positive integers
		\[
		n=\lambda_1+\lambda_2+\cdots+\lambda_r
		\]
		with
		\[
		\lambda_1\ge \lambda_2\ge \cdots \ge \lambda_r\ge 1.
		\]
		The integers \(\lambda_i\) are called the \emph{parts} of the partition.
	\end{definition}
	
	\begin{example}
		The partitions of \(4\) are:
		\[
		4,\qquad 3+1,\qquad 2+2,\qquad 2+1+1,\qquad 1+1+1+1.
		\]
		So the number of partitions of \(4\) is \(5\).
	\end{example}
	
	\begin{definition}
		Let \(p(n)\) denote the number of partitions of \(n\), with \(p(0)=1\).
	\end{definition}
	
	\subsection{Generating function for partitions}
	
	A standard argument shows that
	\[
	\sum_{n=0}^{\infty} p(n) q^n
	=
	\prod_{m=1}^{\infty}\frac{1}{1-q^m}.
	\]
	
	Indeed, for each integer \(m\ge 1\), the factor
	\[
	\frac{1}{1-q^m}=1+q^m+q^{2m}+q^{3m}+\cdots
	\]
	records how many times the part \(m\) is used in a partition. Multiplying over all \(m\) gives the partition generating function.
	
	The Pentagonal Number Theorem concerns the reciprocal product
	\[
	\prod_{m=1}^{\infty}(1-q^m).
	\]
	
	\section{The Pentagonal Numbers}
	
	\subsection{Ordinary pentagonal numbers}
	
	\begin{definition}
		The \(n\)-th pentagonal number is
		\[
		P_n=\frac{n(3n-1)}{2},
		\qquad n\ge 1.
		\]
	\end{definition}
	
	These begin as
	\[
	1,5,12,22,35,\dots
	\]
	
	They arise from a classical polygonal-number formula. For the Pentagonal Number Theorem, however, what appears naturally is a symmetrized version.
	
	\subsection{Generalized pentagonal numbers}
	
	\begin{definition}
		The \emph{generalized pentagonal numbers} are the integers
		\[
		\frac{k(3k-1)}{2},
		\qquad k\in \ZZ.
		\]
	\end{definition}
	
	If we list them for
	\[
	k=0,1,-1,2,-2,3,-3,\dots
	\]
	we get
	\[
	0,1,2,5,7,12,15,22,26,\dots
	\]
	
	Explicitly:
	\[
	\frac{1(3\cdot 1-1)}{2}=1,\qquad
	\frac{(-1)(3(-1)-1)}{2}=2,
	\]
	\[
	\frac{2(6-1)}{2}=5,\qquad
	\frac{(-2)(-6-1)}{2}=7,
	\]
	and so on.
	
	\begin{remark}
		The generalized pentagonal numbers occur in pairs:
		\[
		\frac{k(3k-1)}{2},\qquad \frac{k(3k+1)}{2},
		\]
		for \(k\ge 1\), because
		\[
		\frac{(-k)(3(-k)-1)}{2}=\frac{k(3k+1)}{2}.
		\]
	\end{remark}
	
	Thus the theorem is often written in the form
	\[
	\prod_{m=1}^{\infty}(1-q^m)
	=
	1+\sum_{k=1}^{\infty}(-1)^k
	\left(q^{k(3k-1)/2}+q^{k(3k+1)/2}\right).
	\]
	
	\section{Statement of the Pentagonal Number Theorem}
	
	\begin{theorem}[Pentagonal Number Theorem]
		In the ring of formal power series,
		\[
		\prod_{m=1}^{\infty}(1-q^m)
		=
		\sum_{k=-\infty}^{\infty}(-1)^k q^{k(3k-1)/2}.
		\]
		Equivalently,
		\[
		\prod_{m=1}^{\infty}(1-q^m)
		=
		1+\sum_{k=1}^{\infty}(-1)^k
		\left(q^{k(3k-1)/2}+q^{k(3k+1)/2}\right).
		\]
	\end{theorem}
	
	Expanding the first terms gives
	\[
	\prod_{m=1}^{\infty}(1-q^m)
	=
	1-q-q^2+q^5+q^7-q^{12}-q^{15}+q^{22}+q^{26}-\cdots
	\]
	There are no other powers of \(q\): only the generalized pentagonal numbers appear.
	
	\section{Why This Is Surprising}
	
	The finite partial products
	\[
	(1-q)(1-q^2)(1-q^3)\cdots(1-q^N)
	\]
	expand into many terms, with cancellations. It is far from obvious that in the infinite limit the remaining coefficients are so sparse and so structured.
	
	The theorem says:
	\begin{itemize}
		\item if \(n\) is not a generalized pentagonal number, then the coefficient of \(q^n\) is \(0\);
		\item if
		\[
		n=\frac{k(3k\pm 1)}{2},
		\]
		then the coefficient is \((-1)^k\).
	\end{itemize}
	
	So the whole power series is completely determined by this remarkable pattern.
	
	\section{Combinatorial Meaning of the Product}
	
	Let us interpret
	\[
	\prod_{m=1}^{\infty}(1-q^m)
	\]
	combinatorially.
	
	For each \(m\), the factor
	\[
	1-q^m
	\]
	means: either do not choose the part \(m\), or choose it once and attach a minus sign.
	
	Therefore the coefficient of \(q^n\) in the product is
	\[
	\sum_{\lambda}(-1)^{\ell(\lambda)},
	\]
	where the sum runs over partitions \(\lambda\) of \(n\) into \emph{distinct} parts, and \(\ell(\lambda)\) denotes the number of parts.
	
	So the coefficient of \(q^n\) is:
	\[
	(\text{number of partitions of \(n\) into an even number of distinct parts})
	-
	(\text{number of partitions of \(n\) into an odd number of distinct parts}).
	\]
	
	The Pentagonal Number Theorem therefore says that this signed difference is almost always \(0\), except at the generalized pentagonal numbers.
	
	\section{A First Examples of the Coefficients}
	
	\begin{example}
		For \(n=3\), the partitions into distinct parts are
		\[
		3,\qquad 2+1.
		\]
		So the signed contribution is
		\[
		(-1)^1+(-1)^2=-1+1=0.
		\]
		Indeed, \(3\) is not a generalized pentagonal number.
	\end{example}
	
	\begin{example}
		For \(n=5\), the partitions into distinct parts are
		\[
		5,\qquad 4+1,\qquad 3+2.
		\]
		The signed sum is
		\[
		(-1)^1+(-1)^2+(-1)^2=-1+1+1=1.
		\]
		Since \(5=\frac{2(3\cdot 2-1)}{2}\), the theorem predicts coefficient \(+1\).
	\end{example}
	
	\begin{example}
		For \(n=6\), the partitions into distinct parts are
		\[
		6,\qquad 5+1,\qquad 4+2,\qquad 3+2+1.
		\]
		The signed sum is
		\[
		-1+1+1-1=0.
		\]
		Again this agrees with the theorem.
	\end{example}
	
	\section{A Proof via Euler's Recurrence Argument}
	
	We first explain the theorem in a way closely tied to partitions and generating functions. Then later we give the classic involution proof.
	
	Let
	\[
	F(q)=\prod_{m=1}^{\infty}(1-q^m).
	\]
	Suppose
	\[
	F(q)=\sum_{n=0}^{\infty} a_n q^n.
	\]
	Then
	\[
	\left(\sum_{n=0}^{\infty} p(n)q^n\right)\left(\sum_{n=0}^{\infty} a_n q^n\right)=1,
	\]
	because
	\[
	\sum_{n=0}^{\infty} p(n)q^n=\prod_{m=1}^{\infty}\frac{1}{1-q^m}.
	\]
	
	Thus the coefficients \(a_n\) are determined by the condition that
	\[
	\sum_{k=0}^{n} p(k)a_{n-k}=0
	\qquad (n\ge 1),
	\]
	with \(a_0=1\).
	
	Euler observed experimentally that the \(a_n\) are exactly
	\[
	a_n=
	\begin{cases}
		(-1)^m, & n=\dfrac{m(3m-1)}{2}\text{ or }\dfrac{m(3m+1)}{2},\\[1em]
		0, & \text{otherwise}.
	\end{cases}
	\]
	
	This leads to the partition recurrence
	\[
	p(n)=p(n-1)+p(n-2)-p(n-5)-p(n-7)+p(n-12)+p(n-15)-\cdots
	\]
	where the offsets are generalized pentagonal numbers.
	
	This recurrence is one of the main uses of the theorem.
	
	\section{Euler's Recurrence for the Partition Function}
	
	\begin{corollary}[Euler recurrence]
		For \(n\ge 1\),
		\[
		p(n)
		=
		\sum_{k=1}^{\infty}(-1)^{k-1}
		\left(
		p\!\left(n-\frac{k(3k-1)}{2}\right)
		+
		p\!\left(n-\frac{k(3k+1)}{2}\right)
		\right),
		\]
		where \(p(m)=0\) for \(m<0\).
	\end{corollary}
	
	\begin{proof}
		Starting from
		\[
		\left(\sum_{n=0}^{\infty}p(n)q^n\right)
		\left(1+\sum_{k=1}^{\infty}(-1)^k
		\left(q^{k(3k-1)/2}+q^{k(3k+1)/2}\right)\right)=1,
		\]
		compare coefficients of \(q^n\) for \(n\ge 1\).
	\end{proof}
	
	\begin{example}
		Let us compute \(p(5)\):
		\[
		p(5)=p(4)+p(3)-p(0).
		\]
		Since
		\[
		p(4)=5,\qquad p(3)=3,\qquad p(0)=1,
		\]
		we get
		\[
		p(5)=5+3-1=7.
		\]
	\end{example}
	
	\section{Ferrers Diagrams}
	
	To explain the combinatorial proof, we need diagrams.
	
	\begin{definition}
		The \emph{Ferrers diagram} of a partition
		\[
		\lambda=(\lambda_1,\lambda_2,\dots,\lambda_r)
		\]
		is the array of dots or boxes with \(\lambda_i\) boxes in row \(i\), left-justified.
	\end{definition}
	
	\begin{example}
		The partition
		\[
		7=4+2+1
		\]
		has Ferrers diagram
		\[
		\begin{array}{llll}
			\bullet & \bullet & \bullet & \bullet\\
			\bullet & \bullet\\
			\bullet
		\end{array}
		\]
	\end{example}
	
	For partitions into distinct parts, the row lengths are strictly decreasing.
	
	\section{Durfee Squares}
	
	\begin{definition}
		The \emph{Durfee square} of a partition is the largest square that fits inside its Ferrers diagram.
	\end{definition}
	
	\begin{example}
		For
		\[
		\lambda=(6,4,3,1),
		\]
		the Durfee square has side length \(3\), since the first three rows have lengths at least \(3\), but a \(4\times 4\) square does not fit.
	\end{example}
	
	Durfee squares are useful in the involution proof because they organize partitions according to their geometric shape.
	
	\section{Euler's Involution: Outline}
	
	Let \(D(n)\) be the set of partitions of \(n\) into distinct parts. We want to evaluate
	\[
	\sum_{\lambda\in D(n)}(-1)^{\ell(\lambda)}.
	\]
	
	Euler's idea is to pair off most partitions in \(D(n)\) so that each pair contributes
	\[
	(-1)^r+(-1)^{r+1}=0.
	\]
	Only a small exceptional family remains unmatched, and those exceptional partitions occur exactly when \(n\) is a generalized pentagonal number.
	
	That is the source of the theorem.
	
	\section{Description of the Exceptional Partitions}
	
	The unmatched partitions in Euler's involution turn out to be of the forms
	\[
	(2m-1,\,2m-2,\,\dots,\,m)
	\]
	and
	\[
	(2m,\,2m-1,\,\dots,\,m+1).
	\]
	
	Let us compute their sizes.
	
	For the first:
	\[
	(2m-1)+(2m-2)+\cdots+m
	=
	\sum_{j=m}^{2m-1} j.
	\]
	This is
	\[
	\frac{(m+(2m-1))m}{2}
	=
	\frac{(3m-1)m}{2}.
	\]
	
	For the second:
	\[
	(2m)+(2m-1)+\cdots+(m+1)
	=
	\sum_{j=m+1}^{2m} j
	=
	\frac{((m+1)+2m)m}{2}
	=
	\frac{m(3m+1)}{2}.
	\]
	
	These are exactly the two forms of generalized pentagonal numbers.
	
	\section{Euler's Involution in More Detail}
	
	We now sketch the classic involution. Let
	\[
	\lambda=(\lambda_1>\lambda_2>\cdots>\lambda_r)
	\]
	be a partition into distinct parts.
	
	Define:
	\begin{itemize}
		\item \(s=r\), the number of parts;
		\item \(t=\lambda_r\), the smallest part.
	\end{itemize}
	
	The Ferrers diagram has a diagonal of some length, and one compares the smallest part with the number of parts. Roughly:
	\begin{itemize}
		\item if the smallest part is larger than the number of parts, remove \(1\) from each part and append a new smallest part;
		\item if the smallest part is at most the number of parts, reverse the process by deleting the smallest part and adding \(1\) to several parts.
	\end{itemize}
	
	This operation changes the number of parts by \(1\), hence changes the sign. It is an involution on all non-exceptional partitions.
	
	The exceptional cases are precisely the staircase-like partitions described above, and these yield the generalized pentagonal numbers.
	
	\section{A More Precise Involution Formulation}
	
	Let
	\[
	\lambda=(\lambda_1>\lambda_2>\cdots>\lambda_r)
	\]
	be a partition into distinct parts.
	
	There is a standard map defined on non-exceptional partitions:
	\begin{enumerate}[label=\arabic*.]
		\item If \(\lambda_r>r\), replace
		\[
		(\lambda_1,\lambda_2,\dots,\lambda_r)
		\]
		by
		\[
		(\lambda_1-1,\lambda_2-1,\dots,\lambda_r-1,r).
		\]
		\item If \(\lambda_r\le r\), reverse this process.
	\end{enumerate}
	
	One checks:
	\begin{itemize}
		\item the resulting partition still has distinct parts except in the exceptional staircase cases;
		\item applying the map twice returns to the original partition;
		\item the parity of the number of parts changes.
	\end{itemize}
	
	Hence all non-exceptional partitions cancel in the signed sum.
	
	The only survivors are the two staircase families:
	\[
	(2m-1,2m-2,\dots,m),
	\qquad
	(2m,2m-1,\dots,m+1).
	\]
	These have \(m\) parts, so each contributes sign \((-1)^m\).
	
	Thus
	\[
	\prod_{m=1}^{\infty}(1-q^m)
	=
	1+\sum_{m=1}^{\infty}(-1)^m
	\left(q^{m(3m-1)/2}+q^{m(3m+1)/2}\right).
	\]
	
	This proves the theorem.
	
	\section{Proof of the Pentagonal Number Theorem}
	
	\begin{proof}[Combinatorial proof]
		Expand
		\[
		\prod_{m=1}^{\infty}(1-q^m)
		\]
		as a generating function for partitions into distinct parts, with sign \((-1)^{\ell(\lambda)}\). Thus the coefficient of \(q^n\) is
		\[
		\sum_{\lambda\in D(n)}(-1)^{\ell(\lambda)}.
		\]
		
		Euler's involution pairs all non-exceptional partitions in \(D(n)\) into pairs whose numbers of parts differ by \(1\). Therefore each pair contributes zero to the signed sum.
		
		The only partitions not paired are the two staircase families:
		\[
		(2m-1,2m-2,\dots,m),
		\qquad
		(2m,2m-1,\dots,m+1).
		\]
		These have sizes
		\[
		\frac{m(3m-1)}{2},
		\qquad
		\frac{m(3m+1)}{2},
		\]
		respectively, and each has \(m\) parts, hence contributes sign \((-1)^m\).
		
		Therefore
		\[
		\prod_{m=1}^{\infty}(1-q^m)
		=
		1+\sum_{m=1}^{\infty}(-1)^m
		\left(q^{m(3m-1)/2}+q^{m(3m+1)/2}\right),
		\]
		which is exactly the Pentagonal Number Theorem.
	\end{proof}
	
	\section{First Terms of the Expansion}
	
	Let us write out the first terms:
	\[
	\prod_{m=1}^{\infty}(1-q^m)
	=
	1-q-q^2+q^5+q^7-q^{12}-q^{15}+q^{22}+q^{26}-\cdots
	\]
	
	The exponents are:
	\[
	0,\ 1,\ 2,\ 5,\ 7,\ 12,\ 15,\ 22,\ 26,\dots
	\]
	coming from
	\[
	\frac{k(3k-1)}{2},
	\qquad k\in \ZZ.
	\]
	
	The signs occur in pairs:
	\[
	+,\ -,-,\ +,+,\ -,-,\ +,+,\dots
	\]
	more precisely:
	\[
	(-1)^m \quad \text{for the pair } \frac{m(3m-1)}{2},\ \frac{m(3m+1)}{2}.
	\]
	
	\section{Computing with the Theorem}
	
	\begin{example}
		Using the theorem,
		\[
		\prod_{m=1}^{\infty}(1-q^m)
		=
		1-q-q^2+q^5+q^7-\cdots
		\]
		so the coefficient of \(q^6\) is \(0\).
	\end{example}
	
	\begin{example}
		The coefficient of \(q^{12}\) is \(-1\), because
		\[
		12=\frac{3(3\cdot 3-1)}{2}=\frac{3\cdot 8}{2}=12.
		\]
	\end{example}
	
	\begin{example}
		The coefficient of \(q^{15}\) is also \(-1\), because
		\[
		15=\frac{3(3\cdot 3+1)}{2}.
		\]
	\end{example}
	
	\section{The Theorem as an Identity of Formal Power Series}
	
	It is important to note that the theorem is best interpreted formally. In the ring of formal power series \(\ZZ[[q]]\), the infinite product
	\[
	\prod_{m=1}^{\infty}(1-q^m)
	\]
	is well-defined because the coefficient of \(q^n\) depends only on finitely many factors.
	
	So the theorem is not merely an analytic identity for \(|q|<1\), though it is true analytically there as well. It is an exact algebraic identity in formal power series.
	
	\begin{remark}
		This formal viewpoint is standard in generatingfunctionology and combinatorics.
	\end{remark}
	
	\section{Connection with the Jacobi Triple Product}
	
	The Pentagonal Number Theorem can also be derived from the Jacobi triple product:
	\[
	\sum_{k=-\infty}^{\infty} z^k q^{k^2}
	=
	\prod_{m=1}^{\infty}(1-q^{2m})(1+zq^{2m-1})(1+z^{-1}q^{2m-1}).
	\]
	
	By a suitable specialization, one obtains Euler's identity. This places the Pentagonal Number Theorem in a broader family of theta-function identities.
	
	Although we do not pursue that proof here, it is worth knowing that Euler's theorem sits at the entrance to a vast subject.
	
	\section{Geometric Intuition for the Pentagonal Numbers}
	
	Why pentagonal numbers? The short answer is: not because pentagons are visually built into the product, but because the exceptional staircase partitions have sizes
	\[
	\frac{m(3m\pm 1)}{2},
	\]
	which are exactly the generalized pentagonal numbers.
	
	So the pentagonal numbers arise from the geometry of Ferrers diagrams under Euler's involution, rather than from any immediate polygonal construction on the product side.
	
	\section{A Table of Small Values}
	
	The beginning of
	\[
	\prod_{m=1}^{\infty}(1-q^m)
	\]
	looks like this:
	
	\[
	\begin{array}{c|c}
		n & [q^n]\prod_{m=1}^{\infty}(1-q^m)\\
		\hline
		0 & 1\\
		1 & -1\\
		2 & -1\\
		3 & 0\\
		4 & 0\\
		5 & 1\\
		6 & 0\\
		7 & 1\\
		8 & 0\\
		9 & 0\\
		10 & 0\\
		11 & 0\\
		12 & -1\\
		13 & 0\\
		14 & 0\\
		15 & -1
	\end{array}
	\]
	
	This matches the theorem perfectly.
	
	\section{Applications}
	
	\subsection{Partition recurrence}
	
	The main application is Euler's recurrence for \(p(n)\), already discussed.
	
	\subsection{Congruences and modular ideas}
	
	The infinite products
	\[
	\prod_{m=1}^{\infty}(1-q^m)
	\quad\text{and}\quad
	\prod_{m=1}^{\infty}\frac{1}{1-q^m}
	\]
	are fundamental in modular form theory and in the study of partition congruences such as Ramanujan's:
	\[
	p(5n+4)\equiv 0 \pmod 5,
	\]
	\[
	p(7n+5)\equiv 0 \pmod 7,
	\]
	\[
	p(11n+6)\equiv 0 \pmod{11}.
	\]
	
	\subsection{Combinatorial involutions}
	
	The theorem is a classical example of a \emph{sign-reversing involution}, one of the central techniques in enumerative combinatorics.
	
	\section{Summary}
	
	Let us summarize the main points.
	
	\begin{enumerate}[label=\arabic*.]
		\item The generating function for partitions is
		\[
		\sum_{n=0}^{\infty}p(n)q^n
		=
		\prod_{m=1}^{\infty}\frac{1}{1-q^m}.
		\]
		\item The reciprocal product
		\[
		\prod_{m=1}^{\infty}(1-q^m)
		\]
		is the signed generating function for partitions into distinct parts.
		\item Euler's involution pairs off almost all such partitions, leaving only the exceptional staircase partitions.
		\item The sizes of these exceptional partitions are the generalized pentagonal numbers
		\[
		\frac{k(3k-1)}{2},\qquad k\in \ZZ.
		\]
		\item Hence
		\[
		\prod_{m=1}^{\infty}(1-q^m)
		=
		\sum_{k=-\infty}^{\infty}(-1)^k q^{k(3k-1)/2}.
		\]
	\end{enumerate}
	
	That is the Pentagonal Number Theorem.
	
	\section{Final Theorem Statement}
	
	\begin{theorem}[Euler's Pentagonal Number Theorem]
		In \(\ZZ[[q]]\),
		\[
		\prod_{m=1}^{\infty}(1-q^m)
		=
		\sum_{k=-\infty}^{\infty}(-1)^k q^{k(3k-1)/2}
		=
		1+\sum_{k=1}^{\infty}(-1)^k
		\left(q^{k(3k-1)/2}+q^{k(3k+1)/2}\right).
		\]
	\end{theorem}
	
	\section{Exercises}
	
	\begin{exercise}
		List all partitions of \(6\) into distinct parts and verify that their signed count is \(0\).
	\end{exercise}
	
	\begin{exercise}
		Verify directly that the coefficient of \(q^7\) in
		\[
		\prod_{m=1}^{\infty}(1-q^m)
		\]
		is \(1\).
	\end{exercise}
	
	\begin{exercise}
		Show that the partition
		\[
		(2m-1,2m-2,\dots,m)
		\]
		has size \(m(3m-1)/2\).
	\end{exercise}
	
	\begin{exercise}
		Show that the partition
		\[
		(2m,2m-1,\dots,m+1)
		\]
		has size \(m(3m+1)/2\).
	\end{exercise}
	
	\begin{exercise}
		Use the Pentagonal Number Theorem to compute \(p(6)\) from Euler's recurrence.
	\end{exercise}
	
	\begin{exercise}
		Write out the first ten generalized pentagonal numbers.
	\end{exercise}
	
	\begin{exercise}
		Explain why the Pentagonal Number Theorem is naturally a statement in formal power series.
	\end{exercise}
	
	\section{Selected Answers}
	
	\begin{itemize}
		\item The first generalized pentagonal numbers are
		\[
		0,1,2,5,7,12,15,22,26,35.
		\]
		
		\item For \(m=3\),
		\[
		(5,4,3)
		\]
		has size
		\[
		5+4+3=12=\frac{3(3\cdot 3-1)}{2}.
		\]
		
		\item Also for \(m=3\),
		\[
		(6,5,4)
		\]
		has size
		\[
		6+5+4=15=\frac{3(3\cdot 3+1)}{2}.
		\]
	\end{itemize}
	
	\section*{One-Sentence Takeaway}
	
	Euler's Pentagonal Number Theorem says that the infinite product \(\prod_{m\ge 1}(1-q^m)\), which is the signed generating function for partitions into distinct parts, collapses to a sparse series supported exactly on the generalized pentagonal numbers.
	
\end{document}